\section{The Role of the Agent Harness}

\subsection{Harness Comparison: Leading, Native, and Open-Source Harnesses}
\label{sec:harness_ablation}


\begin{wrapfigure}[17]{r}{0.4\textwidth}
    \centering
    \includegraphics[width=\linewidth]{Figures/harness_ablation.pdf}
    \caption{Coding harness comparison.}
    \label{fig:harness_ablation}
\end{wrapfigure}

To measure the effect of harness choice, we compare three harness settings for Claude-Opus-4.7, GPT-5.5, and Kimi-K2.7-Code: the shared Claude Code harness (v2.1.152), each model's native harness, and model-agnostic open-source OpenCode harness (v1.17.18).
The native harnesses are Claude Code for Opus, Codex CLI (v0.142.4) for GPT, and Kimi Code CLI (v0.24.1) for Kimi.
Across conditions, we hold all 36 tasks, the execution environment, resource limits, and three-rollout protocol fixed.


\textbf{The three harness settings achieve comparable aggregate performance and preserve model rankings, differing mainly in run-to-run stability.}
As shown in Figure~\ref{fig:harness_ablation}, best@3 scores vary little across harnesses: the largest difference for any model is $0.035$.
By contrast, avg@3 is more sensitive to harness choice: relative to Claude Code, the native harness and OpenCode raise it by $0.019$ and $0.014$ for GPT-5.5, and by $0.055$ and $0.046$ for Kimi-K2.7-Code, showing that both model-native and open-source harnesses improve performance stability, particularly for Kimi.
Nevertheless, Opus, GPT, and Kimi retain the same ordering across all three harness settings under both metrics, suggesting that harness choice primarily affects run-to-run stability rather than relative model ordering in this comparison.
Appendix~\ref{app:harness_by_category} reports the category-level results, where the best-performing harness can vary across task types for the same model.


\subsection{Auto Harness}

The harness has recently emerged as a lever for improving agent behavior without retraining the model itself~\citep{smallmodels2026}. A growing line of work explores evolving the harness automatically rather than hand-engineering it~\citep{zhang2026selfharness, vero2026, metaharness2026}. 
Building on our prior work on harness evolution \citep{chen2026coharness,wang2026rethinking}, we preliminarily explore automated harness evolution for long-horizon research tasks.

\begin{wrapfigure}[19]{r}{0.44\textwidth}
    \centering
    \vspace{-0.8\baselineskip}
    \includegraphics[width=\linewidth]{Figures/harness_decay.pdf}
    \caption{Gain of the evolved harness over the original harness across four transfer settings, from the three seed tasks it was evolved on out to unrelated task families. The gain is largest on the seed tasks and still transfers to held-out same-model and cross-model System Optimization tasks, but does not clearly generalize to unrelated task families.}
    \label{fig:auto_harness}
\end{wrapfigure}

We add an outer loop, driven by Claude-Opus-4.8, that automatically optimizes the harness.
Starting from the Claude Code harness running LongCat-2.0, the optimizer inspects agent behavior on three randomly chosen System Optimization tasks and evolves the harness over just four rounds, refining only its preamble, a few standing in-context rules, and a thin layer of hooks.
The resulting harness is generic and task-agnostic, converging on three simple interventions: identify what the verifier actually rewards, attempt one larger structural change when the score plateaus, and protect the best verified state against a late regressing edit.
We then freeze this evolved harness and apply it unchanged to each task for evaluation.

Figure~\ref{fig:auto_harness} reports the gain of the evolved harness over the original harness across four settings.
On the three seed tasks it lifts avg@3 by $+0.12$, and the gain still carries to the remaining same-model System Optimization tasks ($+0.06$ avg@3) and to a different model, GPT-5.5 ($+0.03$ avg@3).
On unrelated task families, however, it no longer generalizes, showing no clear gain: a harness evolved on only three System Optimization tasks captures little of what other families reward, and a broader, more diverse seed set would likely be needed.
Overall, a four-round search already yields gains that transfer across System Optimization tasks and to a new model, pointing to clear headroom for deeper harness optimization.


\subsection{How Harnesses Support the Research Loop}

Trajectory inspection suggests that general-purpose harnesses mainly support the research loop in three ways. \textbf{Tool interaction and failure recovery:} failed commands and invalid tool inputs are returned as explicit observations, allowing the agent to revise its actions and continue the loop. 
\textbf{Context management:} harnesses compress and organize growing interaction histories, helping preserve useful information over long trajectories. \textbf{Research loop management:} explicit task mechanisms help agents decompose complex goals, track progress, and maintain plans across many experiments. Across the 756 Claude Code trajectories, \texttt{TaskCreate} and \texttt{TaskUpdate} were invoked 2,711 and 4,632 times, respectively, while OpenCode and Kimi Code CLI provide lighter todo mechanisms for similar purposes.


Beyond the general harness, the Auto Harness optimization described above produced an evolved harness with controls tailored specifically to auto research.
First, it \textbf{strengthens version control within the research loop.} The harness instructs the agent to save each verified improvement, isolate risky changes in separate commits, and restore unsuccessful experiments, with the final instruction ``Before you finish, restore your best'' preventing late changes from replacing a stronger verified result. Second, it \textbf{helps the agent escape local optima.} After every five new commits, a hook asks the agent to reassess whether progress has plateaued and to attempt a larger structural change when local refinement has stalled.
On \texttt{agent\_tool\_routing}, this reflection was followed by a switch from Python refinement to native C, after which the score increased from approximately $0.37$ to $0.68$.

Overall, the harness stabilizes long research loops and supports task management. 
The Auto Harness further shows that tailoring a harness to the specific needs of auto research can help agents escape local optima, highlighting the potential of specialized harness design.
